\section{研究背景及意义}

\subsection{向量自回归模型的应用背景}

向量自回归（Vector Autoregression, VAR）模型自 Sims（1980）提出以来，已成为宏观经济学与金融实证研究中分析多变量动态交互关系的核心工具。与单方程时间序列模型不同，VAR 模型将系统中每一个变量均视为内生变量，允许所有变量之间存在相互滞后影响，从而能够完整刻画宏观经济变量之间的动态传导机制。一个包含 $N$ 个内生变量且滞后阶数为 $p$ 的标准 VAR 模型，其在时刻 $t$ 的生成机制可以表示为：
\begin{equation}
Y_{t}=c+\Phi_{1}Y_{t-1}+\Phi_{2}Y_{t-2}+\cdots+\Phi_{p}Y_{t-p}+\epsilon_{t},
\end{equation}
其中，$Y_t$ 是维度为 $N \times 1$ 的观测向量，$c$ 是 $N \times 1$ 的常数项向量，$\Phi_i$ 是维度为 $N \times N$ 的滞后系数矩阵，而 $\epsilon_t$ 代表服从均值为零且协方差矩阵为 $\Sigma$ 的残差向量。凭借其灵活的结构设定，VAR 模型在货币政策传导分析、资产价格波动溢出以及宏观预测等领域取得了大量成功的实证应用。

然而，随着信息采集技术与高频金融交易系统的全面普及，实证研究所面对的数据环境发生了根本性转变——研究人员不再局限于三至五个总量指标，而是可以获取包含成百上千个经济变量的大型时间序列数据集。美联储圣路易斯分行维护的 FRED-MD 整合了 128 个标准化美国宏观经济月度序列；覆盖 18 个主要经济体、按购买力平价计算占全球产出 67\% 的全球 EPU 网络、CEIC 与 WIND 等数据库乃至高频金融面板，均将截面变量维度 $N$ 推入了高维乃至超高维的范畴。这一变化使得 VAR 建模面临严峻的参数爆炸困境：模型待估参数总量以 $O(N^2p)$ 的速率递增，当 $N=10$、$p=1$ 时系数矩阵即包含 100 个独立参数，有限样本长度 $T$ 根本无法支撑如此庞大的参数空间，残差自由度大幅萎缩，协方差矩阵行列式趋于零，传统大样本渐近理论下的统计推断完全失效。

为克服这一"维度灾难"，现代高维计量经济学引入了两条主要的降维路线。其一为\textbf{稀疏}（Sparse）约束：基于经济系统存在局部连通性的先验假设，通过 Lasso 惩罚将经济关联微弱的冗余参数压缩至零，在单一优化步骤中同步完成变量筛选与参数估计。其二为\textbf{低秩}（Low-Rank）约束：假定系数矩阵 $\Phi$ 受少数几个宏观共同因子主导，可被严格分解为
\begin{equation}
\Phi = U V^{\top},
\end{equation}
其中 $U$ 为 $N \times r$ 的因子载荷矩阵，$V$ 为 $Np \times r$ 的因子动态矩阵，潜在因子数 $r \ll N$；通过核范数正则化或截断奇异值分解提炼出主导系统演化的低维公共结构。两类方法均能在高维场景下有效压缩参数空间，但也为后续的统计推断带来了新的理论挑战，即正则化操作改变了参数空间的拓扑性质，使得传统渐近理论不再适用。

\subsection{结构性变化检验的必要性}

VAR 模型的有效性高度依赖于一个隐含的基础假设——描述变量间动态关联的参数在整个样本期内保持严格不变。然而，宏观经济系统不可避免地会遭受货币政策框架转型、国际贸易摩擦升级、突发性公共卫生危机或金融监管制度重构等外生冲击，导致 VAR 模型中的系数矩阵 $\Phi$、截距项或残差协方差矩阵 $\Sigma$ 出现跨时期的非连续漂移。这种参数的系统性重置在计量经济学中被定义为"结构性变化"（Structural change）或"结构性断裂"（Structural break）。

若忽视这种参数突变，强行以单一约束模型拟合跨越不同经济机制的历史序列，所得估计值将仅是断裂前后真实参数的某种无意义加权平均，既丧失对特定时期的经济学解释力，又破坏时间序列的平稳性前提，导致拟合残差出现严重的自相关与异方差积聚，样本外预测随之产生灾难性失真。因此，在开展任何高维动态系统实证分析之前，检验并准确定位结构性变化，是确保模型推断有效性的必要步骤。

在时间序列结构性变化检验的理论体系中，基于已知断点（Known breakpoint）的检验方法是最核心的基础框架。这一方法由计量经济学家 Chow（1960）提出，其逻辑严密而直接：当研究者根据宏观历史事件先验地确定断点时刻 $t^*$ 时，将全长为 $T$ 的序列在 $t^*$ 处切分为两段（子样本长度分别为 $T_1 = t^* - p$ 与 $T_2 = T - t^*$），分别拟合参数受约束的全局模型（原假设 $H_0$：参数恒定）与参数不受约束的分段模型（备择假设 $H_1$：系数矩阵在 $t^*$ 前后从 $\Phi_1$ 转变为 $\Phi_2$），进而构造似然比检验统计量：
\begin{equation}
LR = 2\bigl(\ln L_U(t^*) - \ln L_R\bigr),
\end{equation}
其中 $\ln L_U(t^*)$ 为非约束模型的对数似然值，$\ln L_R$ 为约束模型的对数似然值。若 $LR$ 超出预设显著性水平下的临界值，则拒绝参数结构稳定的原假设。

已知断点检验在实证研究中获得了广泛应用。1973 年石油危机被确立为美国潜在生产率增速永久性放缓的断裂点；1989 年美联储平滑利率调整策略的转型通过 Chow 检验得到统计确认；2015 年"811 汇率改革"被植入高维汇率模型以验证在岸与离岸人民币市场溢出结构的变化；2020 年 COVID-19 冲击则被用于剖析全球股市关联网络的断裂与各国央行量化宽松政策的干预效力。

相较于未知断点的 Sup-LR 检验，已知断点方法无需在时间区间 $[\tau_{\min}, \tau_{\max}]$ 内逐一遍历候选断点，彻底规避了多重比较问题，无需对渐近于布朗桥过程的非标准临界值进行额外校正。因此，在相同显著性水平下，已知断点检验具有更高的局部统计功效（Local power），能够敏锐捕捉渐进式改革或温和政策调整所引发的参数微变。

\subsection{高维VAR结构变化检验的挑战与研究意义}

将上述结构性变化检验推广至高维 VAR 场景时，面临两重叠加的技术障碍，使得经典 Chow 检验的推断可靠性受到根本性冲击。

第一重障碍来自\textbf{维度膨胀导致的渐近失真}。随着变量维度 $N$ 上升，基于大样本中心极限定理推导的卡方临界值在有限样本下产生极端失真，典型表现为"过度拒绝"（Over-rejection）——即便数据生成过程在断点前后被严格设定为参数恒定，高维矩阵求逆与行列式计算所积累的伪相关性仍会驱使似然比统计量异常膨胀，使得实际第一类错误率远超名义显著性水平 $5\%$，研究结论中充斥着虚假的结构断裂发现。

第二重障碍来自\textbf{正则化破坏渐近分布}。为克服参数爆炸而施加的 Lasso 稀疏惩罚或截断 SVD 低秩约束，均具有非连续性与非凸的数学特征，根本性地改变了参数空间的拓扑性质。在此非凸空间内，传统的泰勒展开与渐近正态性假设全面失效，检验统计量不再收敛于任何标准分布，所有现成的临界值表在高维正则化模型面前完全失效。

两重障碍的叠加使得高维 VAR 场景下的已知断点检验陷入两难困境：直接套用传统渐近临界值会因维度膨胀而产生大量虚假断裂发现，而现有针对低维模型的推断修正方案又无法处理 Lasso 或低秩约束所引入的非凸参数空间。如何在高维正则化约束下为结构变化检验提供有效且可靠的统计推断，是当前该领域亟待解决的核心理论问题，也构成了本文研究的出发点。
